% =========================================================
% Product Metric Spaces
% Source: volume-iii/analysis/metric-spaces/notes/metric-spaces/
%
% Dependencies:
%   notes-metric-space.tex           — def:metric
%   notes-multiple-metrics.tex       — equivalent metrics
%   notes-convergence.tex            — convergence in metric spaces
% =========================================================

% ---------------------------------------------------------
\subsubsection{Product Metric Spaces}
\label{sec:product-metric}
% ---------------------------------------------------------

% ---- Toolkit ----
\begin{tcolorbox}[colback=gray!6, colframe=gray!40, arc=2pt,
  left=6pt, right=6pt, top=4pt, bottom=4pt,
  title={\small\textbf{Product Metrics --- Quick Reference}},
  fonttitle=\small\bfseries]
\small
\begin{tabular}{l l l l}
\toprule
\textbf{Label} & \textbf{Statement} & \textbf{Proof method} & \textbf{Detail} \\
\midrule
D\ref{def:product-metrics}
  & Three standard metrics on $X \times Y$
  & Definition
  & \hyperref[def:product-metrics]{$\downarrow$} \\[4pt]
P\ref{prop:product-metrics-equiv}
  & $d_1, d_2, d_\infty$ on $X \times Y$ are equivalent
  & Explicit inequalities
  & \hyperref[prop:product-metrics-equiv]{$\downarrow$} \\[4pt]
P\ref{prop:product-convergence}
  & $(x_n, y_n) \to (x,y)$ iff $x_n \to x$ and $y_n \to y$
  & Component-wise
  & \hyperref[prop:product-convergence]{$\downarrow$} \\[4pt]
P\ref{prop:Rn-as-product}
  & $\mathbb{R}^n$ is a metric space via iterated products
  & Induction
  & \hyperref[prop:Rn-as-product]{$\downarrow$} \\[4pt]
\bottomrule
\end{tabular}
\end{tcolorbox}

% ---------------------------------------------------------
\paragraph{Motivation.}
So far every metric space we have encountered has been a single set
with a single metric.
But mathematics constantly asks us to work with pairs:
a point in $\mathbb{R}^2$ is a pair $(x,y)$,
a point in function space is a pair (input, output).
The natural question is: if $(X, d_X)$ and $(Y, d_Y)$ are metric spaces,
how do we measure distance between pairs $(x_1, y_1)$ and $(x_2, y_2)$?

The answer is not unique, and that is a feature, not a bug.
There are several natural choices, and they all produce the same
topology (same convergent sequences, same open sets).
The flexibility lets us pick whichever is most convenient for
the problem at hand.

% ---- Definition: Product metrics ----
\begin{tcolorbox}[colback=propbox, colframe=propborder, arc=2pt,
  left=6pt, right=6pt, top=4pt, bottom=4pt,
  title={\small\textbf{Definition (Product Metrics)}},
  fonttitle=\small\bfseries]
\begin{definition}[Product Metrics]
\label{def:product-metrics}
Let $(X, d_X)$ and $(Y, d_Y)$ be metric spaces.
On the Cartesian product $X \times Y$, define three functions:
\begin{align*}
d_2\bigl((x_1,y_1),(x_2,y_2)\bigr)
  &:= \sqrt{d_X(x_1,x_2)^2 + d_Y(y_1,y_2)^2}
  \quad\text{(Euclidean product metric)},\\[4pt]
d_1\bigl((x_1,y_1),(x_2,y_2)\bigr)
  &:= d_X(x_1,x_2) + d_Y(y_1,y_2)
  \quad\text{(taxicab product metric)},\\[4pt]
d_\infty\bigl((x_1,y_1),(x_2,y_2)\bigr)
  &:= \max\!\bigl(d_X(x_1,x_2),\, d_Y(y_1,y_2)\bigr)
  \quad\text{(supremum product metric)}.
\end{align*}
Each of $d_2$, $d_1$, $d_\infty$ is a metric on $X \times Y$.
\end{definition}
\end{tcolorbox}

\begin{remark*}[Verifying the metric axioms]
Non-negativity, identity of indiscernibles, and symmetry are immediate
from the corresponding properties of $d_X$ and $d_Y$.
The triangle inequality for $d_2$ follows from the Cauchy--Schwarz
inequality applied to the two-component vector $(d_X, d_Y)$.
For $d_1$ and $d_\infty$ the triangle inequality is a direct calculation:
\begin{align*}
d_1((x_1,y_1),(x_3,y_3))
&= d_X(x_1,x_3) + d_Y(y_1,y_3) \\
&\le \bigl(d_X(x_1,x_2) + d_X(x_2,x_3)\bigr) +
     \bigl(d_Y(y_1,y_2) + d_Y(y_2,y_3)\bigr) \\
&= d_1((x_1,y_1),(x_2,y_2)) + d_1((x_2,y_2),(x_3,y_3)).
\end{align*}
The $d_\infty$ case is similar.
\end{remark*}

\begin{remark*}[Intuition for the three metrics]
In $\mathbb{R}^2$ these are the three metrics you already know:
$d_2$ is the straight-line (Euclidean) distance,
$d_1$ is the taxicab distance (walking along a grid),
and $d_\infty$ is the Chebyshev distance (chessboard moves).
The unit balls of each look different --- a circle, a diamond, a square ---
but they all produce the same notion of ``nearby.''
\end{remark*}

% ---- Proposition: equivalence ----

\begin{proposition}[The three product metrics are equivalent]
\label{prop:product-metrics-equiv}
For any $(x_1,y_1), (x_2,y_2) \in X \times Y$:
\[
d_\infty \;\le\; d_2 \;\le\; d_1 \;\le\; 2\, d_\infty.
\]
In particular, $d_1$, $d_2$, $d_\infty$ are equivalent metrics on $X \times Y$.
\end{proposition}

\begin{proof}
Write $a = d_X(x_1,x_2)$ and $b = d_Y(y_1,y_2)$ for brevity.

\medskip
\noindent\textbf{$d_\infty \le d_2$:} \;
$\max(a,b)^2 \le a^2 + b^2$, so $\max(a,b) \le \sqrt{a^2+b^2}$.

\medskip
\noindent\textbf{$d_2 \le d_1$:} \;
$\sqrt{a^2+b^2} \le \sqrt{(a+b)^2} = a+b$
(since $2ab \ge 0$, so $a^2+b^2 \le (a+b)^2$).

\medskip
\noindent\textbf{$d_1 \le 2\,d_\infty$:} \;
$a + b \le \max(a,b) + \max(a,b) = 2\max(a,b)$.

\medskip
These inequalities show that $d_\infty, d_2, d_1$ all induce the same
topology: a sequence is eventually in a $d_\infty$-ball iff it is
eventually in a $d_1$-ball, and so on.
\end{proof}

\begin{remark*}[Consequence: choose whichever is convenient]
Because all three metrics are equivalent, any statement about
convergence, continuity, open/closed sets, or completeness is the
same for all three.
In practice:
\begin{itemize}
  \item $d_\infty$ is often easiest for existence proofs (it controls
        each component independently).
  \item $d_2$ is preferred when geometry matters (angles, inner products).
  \item $d_1$ is natural in probability and measure theory.
\end{itemize}
\end{remark*}

% ---- Proposition: component-wise convergence ----

\begin{proposition}[Convergence in a product is component-wise]
\label{prop:product-convergence}
Let $(x_n, y_n)$ be a sequence in $X \times Y$ (with any of the equivalent
product metrics) and let $(x, y) \in X \times Y$.
Then
\[
(x_n, y_n) \to (x, y)
\;\iff\;
x_n \to x \text{ in } (X, d_X)
\;\text{ and }\;
y_n \to y \text{ in } (Y, d_Y).
\]
\end{proposition}

\begin{proof}
Work with $d_\infty$ for simplicity.

$(\Rightarrow)$: Suppose $(x_n, y_n) \to (x, y)$ in $d_\infty$.
Given $\varepsilon > 0$, choose $N$ so that for $n \ge N$:
$\max(d_X(x_n,x), d_Y(y_n,y)) < \varepsilon$.
Then $d_X(x_n,x) < \varepsilon$ and $d_Y(y_n,y) < \varepsilon$, so both
component sequences converge.

$(\Leftarrow)$: Suppose $x_n \to x$ and $y_n \to y$.
Given $\varepsilon > 0$, choose $N_1$ so $d_X(x_n,x) < \varepsilon$ for
$n \ge N_1$, and $N_2$ so $d_Y(y_n,y) < \varepsilon$ for $n \ge N_2$.
For $n \ge \max(N_1, N_2)$:
$d_\infty((x_n,y_n),(x,y)) = \max(d_X(x_n,x), d_Y(y_n,y)) < \varepsilon$.
\end{proof}

\begin{remark*}[This is not obvious in general topology]
Component-wise convergence is a special feature of product metric spaces.
In the general topological setting (the Tychonoff product), convergence
of a sequence does \emph{not} reduce to component-wise convergence
--- only nets behave correctly.
The product metric is what makes sequences well-behaved here.
\end{remark*}

% ---- Rn as iterated product ----

\begin{proposition}[$\mathbb{R}^n$ as an iterated product]
\label{prop:Rn-as-product}
The Euclidean space $(\mathbb{R}^n, d_2)$ is the iterated product metric
space obtained by taking $\mathbb{R}$ with the standard metric $n$ times:
\[
\mathbb{R}^n = \underbrace{\mathbb{R} \times \mathbb{R} \times \cdots \times \mathbb{R}}_{n}
\]
with the Euclidean product metric
$d_2(x,y) = \sqrt{\sum_{i=1}^n (x_i - y_i)^2}$.
The taxicab metric $d_1(x,y) = \sum_{i=1}^n |x_i - y_i|$
and the supremum metric $d_\infty(x,y) = \max_i |x_i - y_i|$
are equivalent metrics on $\mathbb{R}^n$.
\end{proposition}

\begin{remark*}[Convergence in $\mathbb{R}^n$]
By Proposition~\ref{prop:product-convergence} (applied iteratively),
a sequence of vectors $(x_n^{(k)})$ in $\mathbb{R}^n$ converges iff
each coordinate sequence $(x_n^{(k)})_i$ converges in $\mathbb{R}$.
This reduces all analysis in $\mathbb{R}^n$ to $n$ separate problems
in $\mathbb{R}$ --- one per coordinate.
\end{remark*}

\begin{example}[A sequence in $\mathbb{R}^2$]
\label{ex:product-convergence-R2}
Consider $(x_n, y_n) = (1/n,\, n/(n+1))$ in $\mathbb{R}^2$.

$x_n = 1/n \to 0$ in $\mathbb{R}$.
$y_n = n/(n+1) = 1 - 1/(n+1) \to 1$ in $\mathbb{R}$.

By Proposition~\ref{prop:product-convergence},
$(x_n, y_n) \to (0, 1)$ in $\mathbb{R}^2$
with respect to any of the three product metrics.

Explicitly in $d_2$:
$d_2((x_n,y_n),(0,1)) = \sqrt{(1/n)^2 + (1/(n+1))^2} \to 0$.
\end{example}

\begin{remark*}[Infinite products]
The product construction extends to countably infinite products:
given metric spaces $(X_n, d_n)$ for $n \in \mathbb{N}$, the product
$\prod_{n=1}^\infty X_n$ can be metrized by
\[
d(x, y) = \sum_{n=1}^\infty \frac{1}{2^n} \cdot
           \frac{d_n(x_n, y_n)}{1 + d_n(x_n, y_n)}.
\]
This metric makes convergence component-wise on each $X_n$.
The Hilbert cube $[0,1]^\infty$ is a standard example,
metrized this way, and is a compact metric space.
Infinite products belong to the topology chapter;
we note them here for completeness of the picture.
\end{remark*}
